% Complete corrected elementary-boundary and finite-family induction proofs for Kamp III.2, Lemma 8.

\section{Completion of Kamp's universal-family induction}
\label{sec:annex-p09-kamp-iii2-lemma8-completion}

\begin{annexitem}{annex:p9.kamp.elementary.duality}{thm:p09-kamp-elementary-duality}{Dual construction histories and reversed interval words}
\begin{proof}
An induction on first-order formulas gives
\[
 \mathcal M,s\models\theta^{\mathrm d}
       \quad\Longleftrightarrow\quad
 \mathcal M^{\mathrm{op}},s\models\theta.
 \tag{1}
\]
Unary atoms retain their truth sets. At an order atom, the definition
of the opposite order gives exactly the displayed replacement.
Boolean operations preserve the induction equivalence, and a quantified
variable ranges over the same domain in the two structures, proving
the existential and universal steps. No endpoint or predicate interpretation
is changed in this semantic identity.

We now construct an elementary presentation by induction on a supplied
\(\operatorname{E}\) history. A literal \(\operatorname{At}\)
seed is unchanged. A guarded basic past seed with word \(W\)
becomes the guarded basic future seed with its word reversed; a future
seed becomes the corresponding past seed. To verify this step, reverse
each finite realizing sequence. Its point labels occur in the reversed
order and each open gap is the same set of domain points with the
corresponding label. Reversing a second time gives the converse.
The basic realization theorem therefore proves the required guarded
equivalence with an actual seed of the other basic class.

For a positive-stage limiting past seed, its complete finite realization
inside every interval ending at the current point becomes a reversed
realization inside every interval beginning there. This is the limiting
future seed with the reversed word and the same number of intervals.
The future-to-past step is identical after applying reversal a second
time. At stage zero, the formulas \(\partial^-A\) and
\(\partial^+A\) exchange directly: reverse both strict guards
in their quantified definitions. These checks cover every seed
\(\operatorname{At},\operatorname{Bp},\operatorname{Bf},
\operatorname{Lp},\operatorname{Lf}\) of \(\operatorname{E}\).

At a Boolean construction node, apply the same Boolean constructor to
the already obtained dual presentations. At a proper-substitution node,
apply proper substitution to the two dual presentations. Reversal changes
only order atoms, whereas predicate substitution visits the original unary
atoms and inserts their replacements. Reversing the substituted formula
therefore amounts to reversing the outer formula and every inserted
replacement. This can also be checked at each atomic occurrence using
(1) and the proved substitution identity. Bound names are chosen fresh
on both sides; differences in their choices are alphabetic variants.
The two constructed operands are elementary by induction, so their
proper substitution is an allowed \(\operatorname{E}\) constructor.
This proves the claim for every finite history without adding closure
under arbitrary semantic equivalence to the grammar.

For an elementary interval word, first construct an elementary dual
presentation for each of its labels. In \(\mathcal M\), that
dual label at a point has exactly the original label's truth value
in \(\mathcal M^{\mathrm{op}}\), by (1). Reverse the order
of the word's positions as well. Reversing a realization then verifies
every point and every gap of the resulting labelled word, in both
directions. The elementary-word theorem gives an actual
\(\operatorname{Ep}\) presentation with the reversed endpoints.
Thus any raw dual labels in the display \(W^{\mathrm d}\)
are implemented by their constructed elementary representatives.

A complete linear order remains complete on reversal: a lower bound
in the opposite order is an upper bound in the original, and its
greatest lower bound is the original least upper bound. The other
case exchanges in the same way. Nonemptiness and linearity are also
preserved. Consequently the given \(\mathsf{Com}\) equivalences
in an \(I\) representation apply in the opposite structure and,
by (1), produce the required reversed \(I\) representation.
For basic words this operation changes neither the number of occurrences
nor the number of intervals \cite[ch.~III, pp.~64--68, 92]{Kamp1968}.
\end{proof}
\end{annexitem}

\begin{annexitem}{annex:p9.kamp.unary.accumulation}{thm:p09-kamp-unary-accumulation}{Actual elementary neighborhood and accumulation labels}
\begin{proof}
The one-interval word \((T^K,B,T^K)\) has both elementary
past and elementary future presentations. Its guarded closures give
\(N_B^-\) and \(N_B^+\), since their endpoint truth tests
are tautologies. The closure construction in
\cref{thm:p09-kamp-elementary-intervals} supplies actual
\(\operatorname{E}\) histories. More explicitly, take a
guarded one-interval basic seed on three fresh letters, and properly
substitute \(T^K,B,T^K\) into its three positions. Fresh target
letters occur in none of these replacements, so their successive
substitutions give exactly the indicated formula, up to bound names.
This works when \(B\) is a letter or any supplied elementary label.
In particular \(N_{T^K}^+(w)\) means there is a later point,
and \(N_{T^K}^-(w)\) means there is an earlier one.

Every constituent of \(e_W^+\) and \(e_W^-\) has an
elementary history, and their Boolean combination is allowed in
\(\operatorname{E}\). Suppose first there is no point later
than \(a\). Then the universal guard in \(C_W^+(a)\)
has no instance, so that formula is true; the disjunct
\(\neg N_{T^K}^+(a)\) is true as well.
If there is a later point, choose any one \(b>a\).
The proof of \cref{thm:p09-kamp-fixed-endpoint-limits} establishes
\[
 C_W^+(a)\longleftrightarrow
 A_0(a)\land\lambda_W^+(a)\land
 \exists u\,(a<u<b\land\forall t\,(a<t<u\to B_1(t))).
 \tag{1}
\]
The right side implies \(N_{B_1}^+(a)\).
Conversely, \(A_0(a)\), \(\lambda_W^+(a)\), and a
neighborhood witness \(u>a\) suffice to construct all instances
of \(C_W^+(a)\): for any \(v>a\), take the limiting
suffix strictly below the smaller of \(u,v\), and prepend
\(a\). Its first gap lies in the given \(B_1\) neighborhood.
For \(n=1\) the suffix is a single \(A_1\)-point; for
\(n>1\) it is the complete prescribed suffix realization.
The construction and all gap checks are those of the cited proof,
and require no bound involving the auxiliary choice \(b\).
This proves equivalence with \(e_W^+\) at every point with a future.

At a point with no past, \(C_W^-\) and the no-past disjunct
are both true. If \(C_W^-(b)\) holds at a point with a past,
choose an earlier point \(a<b\) and apply the second fixed-endpoint
limit equivalence. It supplies \(A_n(b)\), \(\lambda_W^-(b)\)
and a past \(B_n\) neighborhood, which give the positive disjunct
of \(e_W^-\). Conversely, from that positive disjunct append the fixed final point
to a limiting prefix lying above the larger of the given past
neighborhood witness and the tested lower bound. The last gap is
contained in the \(B_n\) neighborhood; its endpoint satisfies
\(A_n\). The same \(n=1\) versus \(n>1\) distinction
uses one point or the whole prefix. This is the complete reverse
construction in \cref{thm:p09-kamp-fixed-endpoint-limits} and
proves the formula for \(e_W^-\). Thus the labels work globally,
including both endpoint-vacuity cases.
\end{proof}
\end{annexitem}

\begin{annexitem}{annex:p9.kamp.remote.limits}{thm:p09-kamp-remote-limits}{A least cofinal position and the whole stretched interval}
\begin{proof}
Fix a linear structure, a basic word
\(W=(A_0,B_1,A_1,\ldots,B_n,A_n)\), and \(a<b\).
Let \(\mathcal R\) be the set of all realizing sequences
\[
                a=a_0<a_1<\cdots<a_n<b
\]
of this word. The last point varies, but the first point is fixed.
The condition \(B_W^+(a,b)\) says that for every \(v<b\)
some sequence in \(\mathcal R\) has \(a_n>v\).
In particular its instance \(v=a\) makes \(\mathcal R\)
nonempty. The equivalence with this sequence condition follows
from the proved basic-word realization theorem.

For each \(j=0,\ldots,n\), consider the property
\[
 I_j:\quad\forall v<b\ \exists\mathbf a\in\mathcal R
                                      \text{ such that }a_j>v.
\]
The assumption gives \(I_n\). Property \(I_0\) is false:
its instance \(v=a\) would require \(a_0=a>a\).
Choose the least \(k\in\{1,\ldots,n\}\) for which
\(I_k\) holds. Failure of \(I_{k-1}\) supplies a point
\(r<b\) such that every sequence in \(\mathcal R\)
satisfies \(a_{k-1}\leq r\). Since \(\mathcal R\)
is nonempty and \(a_{k-1}\geq a\) in every sequence,
we also have \(a\leq r\).

Use \(I_k\) at this \(r\) to choose one realization
\(\mathbf a\) with \(d=a_k>r\), and put \(c=a_{k-1}\).
Keep its prefix \((a_0,\ldots,a_{k-1})\).
The old gap \((c,d)\) satisfies \(B_k\), and \(c\leq r<d<b\).
We claim that the entire open interval \((c,b)\) satisfies \(B_k\).
For any \(s\) in that interval, if \(s<d\), its truth
comes from this selected old gap. If \(d\leq s<b\), use
\(I_k\) again at \(v=s\) to choose another realization
\(\mathbf a'\) with \(a'_k>s\). Its preceding point
satisfies \(a'_{k-1}\leq r<d\leq s\), so \(s\)
lies strictly inside its \(k\)-th gap. That gap supplies
\(B_k(s)\). This includes the old endpoint \(s=d\),
which was not tested by the selected gap itself. It proves the
claim without making any intermediate points of different
realizations coincide.

For every \(v<b\), property \(I_k\) supplies a realization
whose entire suffix \((a_k,\ldots,a_n)\) lies inside
\((v,b)\). If \(k<n\), this is the exact limiting-past
condition for the suffix word beginning at \(A_k\).
If \(k=n\), it is the strict unary limit
\(\partial^-A_n(b)\). Thus \(p_{W,k}^-(b)\) holds.
The retained prefix, the stretched \(B_k\) gap and this
terminal label form a realization of \(V_{W,k}^+\) from
\(a\) to \(b\). When \(k=1\) the prefix consists
only of \(a\); when \(k=n\) the terminal label is the
one-letter limit. All candidate indices lie in the stated range.

Conversely, suppose \(V_{W,k}^+\) has a realization from
\(a\) to \(b\). It gives a prefix realizing the original
word up through \(A_{k-1}\), with final prefix point \(c<b\),
the condition \(B_k\) throughout \((c,b)\), and
\(p_{W,k}^-(b)\). For any \(v<b\), let \(q\) be
the larger of \(v,c\). Then \(q<b\). The limiting
label supplies the entire suffix strictly inside \((q,b)\),
or one \(A_n\)-point there if \(k=n\). Append this
suffix to the prefix. The new joining gap lies inside \((c,b)\)
and satisfies \(B_k\); all other point and gap predicates
are retained. The resulting full realization ends above \(q\geq v\)
and below \(b\). This proves every required instance of
\(B_W^+(a,b)\).

Each candidate word uses elementary labels with their declared
limiting histories. Thus it has an actual \(\operatorname{Ep}\)
presentation. There are exactly the \(n\) listed candidates,
and the forward and reverse arguments prove the first equivalence.

For the second equivalence reverse the order and the basic word.
The original condition that full realizations end at \(b\)
and begin arbitrarily close above \(a\) becomes the first
condition for the reversed word from \(b\) toward \(a\).
The first equivalence just proved supplies its candidates in the
opposite order. A suffix of the reversed word is a prefix of
the original word; its limiting-past label becomes the corresponding
limiting-future prefix label by
\cref{thm:p09-kamp-elementary-duality}. If the reversed cut
index is \(j\), the original index is \(k=n-j+1\).
Reversing the candidate gives exactly
\((p_{W,k}^+,B_k,A_k,B_{k+1},\ldots,B_n,A_n)\).
The index map is a bijection of \(\{1,\ldots,n\}\),
so every candidate in the second display occurs once. Reversing
back proves both implications in the original structure.

The source's first-orientation cofinal formula omits the explicit
guard \(a<t\). Under \(a<b\), this makes no difference:
given any \(v<b\), apply its universal condition at the
larger of \(v,a\), which is still below \(b\). Its
witness is then above both \(a\) and \(v\), supplying
our guarded condition. The reverse implication is immediate.
The dual verification uses the smaller of \(v,b\).
On adjacent outer endpoints our guarded condition fails at
\(v=a\), or \(v=b\) in the dual case. No empty-interval
universal premise has been mistaken for a cofinal realization.
Every construction works over \(\mathsf{Lin}\), so its
restriction gives the stated \(I\) representations
\cite[ch.~III, pp.~85, 89--90]{Kamp1968}.
\end{proof}
\end{annexitem}

\begin{annexitem}{annex:p9.kamp.interior.approach}{thm:p09-kamp-interior-approach}{Finite cut choices for a limit beyond an interior boundary}
\begin{proof}
Fix \(a<u\). If no point lies above \(u\), both the
universal condition \(A_W^+(a,u)\) and the no-future disjunct
are true. Suppose instead that some point lies above \(u\).
For each cut of the original word, the prefix condition \(L_c(a,u)\)
has one fixed truth value. Let \(J\) contain exactly its true cuts.

Assume \(A_W^+(a,u)\). At any chosen later bound it supplies
some full realization from \(a\) to a point \(t>u\).
Factor that word at \(u\); a true prefix cut is obtained,
so \(J\ne\varnothing\). If every \(c\in J\) failed
the fixed-start accumulation \(C_{W_c^R}^+(u)\), there
would be a point \(v_c>u\) for which no suffix realization
from \(u\) ended inside \((u,v_c)\). Since \(J\) is
finite and nonempty, its finitely many \(v_c\) have a minimum
\(v>u\). The original approach condition at \(v\) gives
\(u<t<v\) with the full pattern from \(a\) to \(t\).
Its cut factorization supplies some \(c\in J\) and
\(R_c(u,t)\). But \(t<v\leq v_c\), contradicting
that cut's chosen failure. Thus at least one true prefix has
\(C_{W_c^R}^+(u)\), equivalently the actual elementary
label \(e_{W_c^R}^+(u)\).

Conversely, if \(L_c(a,u)\) and that label hold, then for
every \(v>u\) its accumulation condition supplies
\(u<t<v\) with \(R_c(u,t)\). Joining the two factors
gives the full word from \(a\) to \(t\), by
\cref{thm:p09-kamp-cut-factorization}. This proves
\(A_W^+(a,u)\). The argument is the finite-union limit
principle with its minimum-bound proof supplied explicitly.

Every cut prefix is basic, hence elementary. The label at \(u\)
has an elementary history by
\cref{thm:p09-kamp-unary-accumulation}; its insertion at the
right endpoint has property \(I(a,u)\) by the endpoint
substitution lemma. Conjunction and finite disjunction preserve
the supplied interval representations through fusion and concatenation
of their disjunct lists. The no-future label is also elementary
and can be inserted at that same endpoint. Hence the whole displayed
formula has an actual \(I(a,u)\) representation.

Order reversal gives the fixed-final-endpoint version. Each prefix
becomes a suffix, its endpoint label is replaced by its constructed
dual elementary label, and a minimum of finitely many later bounds
becomes a maximum of earlier bounds. The proved duality transports
both directions and the actual interval representation. This covers
the left-approach version with its no-past case as well.
\end{proof}
\end{annexitem}

\begin{annexitem}{annex:p9.kamp.segmented.assembly}{thm:p09-kamp-segmented-assembly}{All interval edges, every interior cut and complete assembly}
\begin{proof}
Start with the supplied finite \(I(x_i,x_j)\) representations.
Their equivalences hold under their endpoint guards over
\(\mathsf{Com}\). The global strict sequence of named points
supplies every required guard. Expand all these finite disjunctions
and distribute their conjunction, retaining every tuple of choices.
If an input representation is empty, use the explicit false elementary
word from the interval-property convention; this retains an equivalent
false branch. Each resulting branch now has only elementary point
conditions and individual elementary interval words on named edges.

We first extend the proved basic cut construction to an elementary
word. Give each occurrence position a distinct fresh unary letter and
form the basic word on those letters. In an arbitrary fixed structure,
interpret each fresh letter by the truth set of its original elementary
label, evaluated in the whole original structure. The basic cut
factorization applies to this interpreted word at every strict interior
named point. Transfer it back by the simultaneous truth-set substitution,
or by a finite sequence of proper substitutions with targets fresh from
all original labels. The output factors are exactly the elementary
words formed by the point-cut and gap-cut operations. In a gap cut,
the repeated interval label supplies its truth at the newly introduced
endpoint and on both adjoining gaps; when the two factors are joined,
that point's truth supplies the interior test at the removed cut point.
In a point cut the point label is retained in both factors. These are
all cases of the proved basic factorization, now with their labels'
actual whole-structure truth sets. Every factor has an elementary
presentation, and no quantified variable inside a label is restricted
to the interval being cut.

For an edge \((x_i,x_j)\), cut at \(x_{i+1}\), then cut
its right-hand factors at \(x_{i+2}\), continuing until all
factors are on adjacent pairs. At each step every factor has finitely
many cuts, and the number of remaining named points inside that edge
falls by one. Retain every sequence of choices and distribute all
resulting finite disjunctions. The two-direction factorization at each
step proves equivalence of the original edge with the final disjunction
of adjacent-edge conjunctions. Repeating this for every original edge
terminates because the original list and all vertex lists are finite.

For each final branch, conjoin all labels imposed at each named point,
using \(T^K\) when no label was imposed. These are elementary
labels by finite conjunction. On each adjacent pair, combine all its
elementary words by the complete fusion construction of Lemma~4.
If the pair has no edge condition, give it the word \((T^K,T^K,T^K)\),
which imposes no additional condition under its already strict guard.
Distribute again over every fusion choice. Insert the outer endpoint
labels into the first and last point positions of the end intervals.
The fusion construction and these endpoint insertions preserve exactly
the conditions on that pair.

Finally concatenate the adjacent words from left to right, using the
combined elementary label at each internal named point in the midpoint
operation of \cref{thm:p09-kamp-elementary-midpoint}. That theorem
proves both directions at each step: a choice of all internal named
points gives the concatenated realization, and a realization of the
concatenated word recovers those points at the designated positions.
The middle points are strictly ordered because each constituent word
has at least one interval. The resulting word has only \(x_0,x_m\)
as free endpoints and has an actual \(\operatorname{Ep}(x_0,x_m)\)
presentation. At \(m=1\), this last stage simply leaves the single
adjacent word and its endpoint labels, without an interior quantifier.

All finite alternatives have been retained through expansion, cutting,
fusion and concatenation. The elementary transformations themselves
are valid over \(\mathsf{Lin}\), but the original \(I\)
representations were used only over \(\mathsf{Com}\).
Thus the final finite disjunction gives precisely the claimed
\(I(x_0,x_m)\) representation. This argument supplies a
general finite assembly form of Lemmas~4--5; it does not promote
the input complete-order equivalences to all linear orders.
\end{proof}
\end{annexitem}

\begin{annexitem}{annex:p9.kamp.localization}{thm:p09-kamp-localization}{Two-sided localization and justified smaller-family calls}
\begin{proof}
Let \(m=M(F)>1\), let \(\phi\) be the removed maximal
member, and fix a complete linear structure with
\(a\leq r<s\leq b\). Each equality case is implemented
by using the same endpoint name, so no equality atom is added.
The remaining members have exactly the three forms specified in
the source family. Their finite enumeration includes each original
member once.

For a left basic member \(\rho(a,t)\), if \(r=a\),
retain that word as a local \(\operatorname{Bp}(r,t)\)
member and add no fixed-side condition. If \(a<r\),
then for every tested \(r<t<s\) the basic cut factorization
at \(a<r<t\) applies. At the fixed assignment to \(a,r\),
select all its prefix cuts that are true. The selected prefixes
give the fixed-side conjunction on \((a,r)\); their matching
suffix cuts give a residual disjunction on \((r,t)\).
Enumerate all possible subsets syntactically. The proof of the
finite-family cut theorem shows, pointwise at this fixed cut,
that the original disjunct can be replaced in the universal matrix
by its selected residuals, provided the corresponding fixed
prefixes are kept in the same alternative.

For a right basic member \(\rho(t,b)\), if \(s=b\),
retain it as a local \(\operatorname{Bp}(t,s)\) member.
If \(s<b\), use the factorization at \(t<s<b\).
Now select the true suffix cuts at \((s,b)\), retain all
of them as fixed-side conjuncts, and use their matching prefix
cuts as residuals on \((t,s)\). This is the other orientation
of the same finite construction. Take the full finite product
of the subset choices for both kinds of crossing inputs.
Point-only \(\operatorname{At}\) members are unchanged,
apart from capture-avoiding renaming of their current-point variable.

For completeness, fix the actual tuple of all true fixed-side cuts
in the current structure. If the original universal condition holds
and a tested point does not satisfy its added label \(\eta\)
or one of the retained native disjuncts, some crossing original
word holds there. Its factorization supplies a residual with a
true fixed factor, so that residual occurs in the chosen tuple.
Thus the tuple's local universal condition holds. Conversely, in
any tuple alternative, a true residual and its retained fixed
factor reconstruct the corresponding original word. Native disjuncts
and \(\eta\) already imply the original matrix directly.
Every tested point therefore satisfies the original universal
condition. If all selected residuals are absent, their disjunction
is false and these same implications rely on the other disjuncts.
This proves equivalence with the complete disjunction of packages
in all four endpoint-placement cases, including no crossing input.
Neither endpoint originally occurring in an input formula has
been silently replaced by a cut variable.

Consider one package before inserting \(\eta\). Its local
family \(G\) consists of the retained At formulas, the native
basic words, the selected basic residual words at the new local
endpoints, and one fresh atom \(Q(t)\). Each native word is
one unchanged output of its originating member. Each crossed
member contributes no longer residual and at most one residual
of its own length, by the exact cut bound. The removed member
\(\phi\) contributes none. Therefore every local word has
length at most \(m\), and at most \(c(F)-1\) members
have length \(m\). Identifying equal formulas in the local
set only decreases this number. The added atom has length one,
strictly below \(m\). The exact source comparison thus gives
\(G\prec F\), by \cref{thm:p09-kamp-family-rank}.

This is an actual family in
\(\operatorname{At}(t)\cup\operatorname{Bp}(r,t)\cup
\operatorname{Bp}(t,s)\), with its three variable names distinct.
No elementary label with hidden endpoint parameters is being
counted as At. The stated induction hypothesis supplies a finite
\(\operatorname{Ep}(r,s)\) representation for its universal
formula. Properly substitute the given one-variable elementary
formula \(\eta\) for the fresh letter \(Q\).
The interval-substitution theorem constructs the resulting
elementary disjuncts and proves their \(\mathsf{Com}\)
equivalence. Other local members and all fixed-side conditions
are unchanged because \(Q\) was fresh from them and from
\(\eta\). The output is exactly the local universal formula
with the added disjunct \(\eta\).

This justifies every recursive call and every subsequent substitution.
The fixed-side conjunctions are retained as interval edges with
their actual endpoints; they are not part of the local family's
rank count and have not been dropped. The construction requires
only finite word operations, the smaller-family representations,
and the already proved proper-substitution closure.
\end{proof}
\end{annexitem}

\begin{annexitem}{annex:p9.kamp.positive.boundaries}{thm:p09-kamp-positive-boundaries}{All eleven supported cases and the preserved brackets}
\begin{proof}
Fix the left basic word \(W\), ordered endpoints \(a<b\),
and the original remaining disjunction \(D\). Write
\(H(t)=\mathcal B_p(W;a,t)\) and let \(P(t)\) be
its unrestricted guarded past closure. The completed foundation
gives the full truth-set descriptors \(\Gamma_i\), their
exhaustion over complete orders, and their regional tests.
We prove the new branch replacement in both directions.

Suppose first that the target universal formula holds in a complete
linear structure. If its distinguished truth set is empty, all
points satisfy \(D\); hence \(Y_1\) holds. If it is
the singleton \(\{c\}\), then \(a<c<b\), \(H(c)\)
holds and all other tested points satisfy \(D\). These are
exactly the conditions of \(Y_5\).

In any of the other nine cases, use the actual unique infimum
\(l\) and supremum \(h\) from its descriptor. They satisfy
\(a\leq l<h\leq b\). A lower type-0 descriptor makes
\(H\)-points coinitial above \(a\), so \(C_W^+(a)\)
holds, equivalently \(e_W^+(a)\). A lower type-1 descriptor
supplies \(H(l)\). A lower type-2 descriptor makes
\(H\)-points occur inside every interval immediately to the
right of \(l\); its explicit quantified formula supplies
\(A_W^+(a,l)\). Thus every lower descriptor implies its
positive lower support, with the original initial endpoint fixed.

At an upper type-0 boundary, the truth set is cofinal below \(b\),
giving the nonempty remote formula \(B_W^+(a,b)\).
Upper type 1 supplies \(H(h)\). Upper type 2 has cofinally
many \(H\)-points below its unattained supremum \(h\),
giving \(B_W^+(a,h)\). These implications follow directly
from the complete first-order descriptor formulas; in particular
the upper type-2 points lie strictly below \(h\), since
the supremum is not attained. Every positive upper support holds.
The foundation's branch equivalence supplies all outer \(D\)
tests, middle \(D\lor P\) tests and the \(D\) tests
at its unattained interior boundaries. These are retained verbatim
in \(Y_i\), so the corresponding positive branch is true.
All nine type pairs have therefore been included.

For the converse, \(Y_1\) implies the target immediately,
since it supplies \(D\) everywhere. In \(Y_5\), the
chosen point is covered by \(H\) and every other point by
the two \(D\) intervals, proving the target.
Now take any true two-boundary branch and a point \(l<t<h\).
Each of the three lower supports supplies a strictly earlier
fixed-start realization: lower type 0 applies \(C_W^+(a)\)
at the bound \(t\), obtaining \(a<q<t\) with \(H(q)\);
lower type 1 uses its point \(q=l\), where \(a<l<t\);
lower type 2 applies \(A_W^+(a,l)\) at \(t\), obtaining
\(l<q<t\) with \(H(q)\).

Every upper support supplies a strictly later one. Upper type 0
applies \(B_W^+(a,b)\) at \(t<b\), obtaining
\(t<q'<b\) with \(H(q')\). Upper type 1 takes
\(q'=h\). Upper type 2 applies \(B_W^+(a,h)\) at
\(t<h\), obtaining \(t<q'<h\) with \(H(q')\).
In all nine combinations there are therefore four strictly ordered
points \(a<q<t<q'\), with the two original fixed-start
realization premises of Lemma~6. That lemma gives
\(H(t)\longleftrightarrow P(t)\).
Thus the middle \(D\lor P\) test implies \(D\lor H\)
at every middle point.

Outside the middle interval, the retained regional matrices impose
\(D\), which suffices irrespective of whether additional
\(H\)-points occur there. An interior boundary of type 1
has the explicit \(H\) witness. A type-2 boundary is tested
by \(D\), so no assertion about \(H\)'s falsity is needed
in this direction. A type-0 boundary is an outer endpoint and
is not tested by the original universal formula. These statements
cover both boundaries in every case: rows 2--4, 6--8 and 9--11
give the three lower supports, while columns represented by rows
2/6/9, 3/7/10 and 4/8/11 give the three upper supports.
They cover exactly the same eleven regional arrangements as the
foundation, including both endpoint tests in case 11.

This proves that each branch individually implies the target over
\(\mathsf{Lin}\). The forward implication used completeness
only to choose a full descriptor and prove that it implies the
positive supports. No required bracketing premise has been deleted:
every such premise was derived explicitly in the converse.
The removal of absence assertions is justified by the retained
\(D\) conditions, not by treating a source formula as equivalent
to a weaker conjunct in isolation.

Finally, each listed support has the claimed elementary type.
The type-0 lower label is in \(\operatorname{E}(a)\).
Type-1 supports are basic interval formulas. Type-2 lower support
has the actual \(I(a,l)\) representation constructed by
the interior-approach theorem; the point has a later point because
\(l<h\). The remote upper supports have the explicit
\(I(a,h)\) or \(I(a,b)\) representations from the remote-limit
theorem. These are established constructions independent of the
full Lemma~8 conclusion. This completes both the semantic branch
replacement and the verification of its support types.
\end{proof}
\end{annexitem}

\begin{annexitem}{annex:p9.kamp.iii2.lemma8}{thm:p09-kamp-iii2-lemma8}{The full smaller-family induction in both orientations}
\begin{proof}
We perform well-founded induction on the exact comparison \(\prec\),
simultaneously for every choice of three distinct variable names and
every finite family in the source's three stated classes. Its
well-foundedness was proved in \cref{thm:p09-kamp-family-rank}.
The induction supplies actual finite elementary representations,
not only assertions that some representation exists.

For the empty family or a point-only family, use
\cref{thm:p09-kamp-family-normalization}. Its explicit word
\((T^K,D_F,T^K)\) proves the base conclusion, including the
empty open-interval condition of the empty family. Now assume
\(M(F)>1\) and that representations have been constructed
for all \(G\prec F\). Choose one maximal basic member
\(\phi\), so deleting it removes one of the \(c(F)\)
maximal occurrences in the finite set.

First suppose \(\phi\in\operatorname{Bp}(x,z)\).
Let \(W\) be its exact word, \(D=D_{F\setminus\{\phi\}}\),
and use the eleven positive branches just proved. The unrestricted
past closure \(P(t)\) has an actual \(\operatorname{E}(t)\)
presentation: it is the guarded basic past seed for this same
word with its bound initial endpoint renamed fresh. This membership
uses the declared basic class, and its expressibility was already
proved in Lemma~1. It does not use the present induction.

Consider each regional universal formula in a branch. Its outer
region is tested with \(D\), and its middle region with
\(D\lor P\). Regions with equal endpoints have no points
and their universal formulas are tautologies; omit them with this
explicit equivalence. Every remaining region has strict endpoints
\(r<s\), each drawn from the named points \(a,l,h,b\),
or from \(a,c,b\) in the singleton case. Apply the localization
theorem with \(\eta=F^K\) in an outer region and the
actual elementary presentation of \(P\) in a middle region.
For each finite tuple it constructs a local family \(G\prec F\).
The induction hypothesis is applicable to that exact family and
its local endpoints; the subsequent fresh-atom substitution is
justified by the interval-substitution theorem. Hence every regional
formula is replaced by a finite disjunction of elementary local
conditions and its retained fixed-side basic interval conditions.

There is no hidden same-rank call. For each region, the selected
maximal \(\phi\) has been removed before localization. Every
other input basic formula contributes at most one residual of its
own length and none longer; each native member is one unchanged
output. At most \(c(F)-1\) local members have length \(M(F)\).
The auxiliary fresh atom has length one, and any coincidences of
output formulas reduce rather than increase the count. Thus the
actual source comparison decreases at every call, even when some
surviving residual has the same length as the original maximal word.
The elementary \(P\) is inserted after the call and is never
miscounted as At. Fixed-side conjunctions stay in the package and
are proved elementary directly, rather than erased from the formula.

Replace every positive boundary support by its established elementary
point or interval representation. At a type-2 interior boundary,
the required point condition \(D(l)\) or \(D(h)\) is
the finite disjunction of the original remaining formulas at that
point. Expand it completely. Each alternative is either an
\(\operatorname{At}\), hence elementary, point label; a
left basic interval from \(a\) to that point; or a right
basic interval from that point to \(b\). Every endpoint in
these basic formulas is retained. In particular a condition on
\((l,b)\) is not renamed to one on \((l,h)\).

Distribute all the finite choices of boundary supports, point
disjuncts and regional localization packages. In every resulting
conjunction, all conditions now have one of two exact forms:
an elementary label at a named point, or an interval formula with
a supplied elementary representation on an edge joining two of
the named points. Some edges can cross other named interior points;
they need not be adjacent. The strict vertex order is already a
conjunct of the branch. Apply
\cref{thm:p09-kamp-segmented-assembly}, which cuts all such edges,
retains every alternative, fuses adjacent conditions and concatenates
them. It returns actual \(\operatorname{Ep}(a,b)\) disjuncts.
For branch 1 there are just the two outer points; for branch 5
there is its singleton witness; for the other nine branches there
are zero, one or two distinct interior bounds according to their
types. The assembly theorem explicitly includes all these sizes.

Each of the eleven branches therefore has a constructed
\(I(a,b)\) representation. Concatenate their finite disjunct
lists. The positive-boundary theorem proves that their disjunction
is equivalent to the original universal formula on every complete
linear structure under \(a<b\). This completes the left-oriented
induction step, with all nine two-boundary combinations, the empty
truth-set branch and the singleton branch included.

For a right-oriented maximal member \(\phi(z,y)\), perform
the construction in the opposite order with the outer endpoints
reversed. Replace each basic input by the basic pattern with its
word reversed and corresponding endpoint order; keep each At
input unchanged. Reversing a realization proves equivalence for
each input under its actual interval guard, so the entire guarded
universal formula is transformed correctly. The distinguished
right member becomes a distinguished left member of the same
maximal occurrence length. Other left and right members exchange
their roles. Thus the transformed family has exactly the same
allowed source forms relative to the reversed endpoints.

Every transformed input has its original length. If forming the
finite set identifies equal transformed formulas, the number at
the maximum can only decrease; at least one maximal-length image
remains. Hence the transformed family's maximum equals \(M(F)\)
and its maximal count is at most \(c(F)\). In its left-oriented
construction, deleting the distinguished maximal image and applying
the local residual bound produces a family with maximal count at
most \(c(F)-1\), or with smaller maximum. Every recursive
call is therefore strictly below the original \(F\), even
if some images have coincided. Repeated identical disjuncts have
the same truth value, so passing to their finite set changes no
guarded universal formula.

The left-oriented construction already proved is a construction from
these smaller-family hypotheses, so it applies directly to this
transformed input. It is not a recursive invocation of the full
theorem at the original rank. The opposite order is complete,
and the induction hypothesis supplies the required smaller-family
representations there. Reverse every resulting elementary word
and construct every dual elementary label using
\cref{thm:p09-kamp-elementary-duality}. This yields an actual
\(\operatorname{Ep}(x,y)\) disjunction in the original order
equivalent to the original right-oriented universal formula.
Thus the second orientation is established with the same source
length comparison and without changing the order class.

Every local step has finite output: the cut sets, products of their
subsets, finite disjunction expansions, edge cuts and fusions are
finite. Every recursive call strictly decreases the well-founded
pair \((M,c)\). To see that the whole construction terminates,
use well-founded induction itself: assuming all smaller calls
terminate, this node makes finitely many such calls and then finitely
many assembly operations, so it terminates. The explicit atomic
base starts this argument. No cardinality of model truth sets or
truth decision is used to run the construction.

The source's distinct-occurrence and auxiliary-atom reductions are
available with their complete correction proof in the foundation.
The present construction also works directly for repeated input
letters, because every cut tracks occurrence positions and every
fresh placeholder is chosen outside the full input. It preserves
all source family members and supplies the corrected finite bracketing
construction for the omitted-case steps. Completeness was used
in the boundary exhaustion; each given \(I\) equivalence was
kept at its \(\mathsf{Com}\) scope. No conclusion of Lemma~9,
Lemma~10 or Theorem~III.2 has been invoked.

Finally, for the one-interval example the sole candidate is
\((A,B,\partial^-C)\). On \(\mathbb R\) with the displayed
interpretations, every \(v<1\) has some \(t\) strictly
between \(\max(v,0)\) and 1, and such \(t\) satisfies
\(C\). The interval predicate holds everywhere, so
\(\phi(0,t)\) holds and the remote condition is true,
although \(C(1)\) is false. On the two-point order the
instance \(v=0\) has no interior witness, and
\(\partial^-C(1)\) fails at its lower bound 0 as well.
This verifies both stated boundaries of that example and completes
the full proof of Lemma~8 \cite[ch.~III, pp.~83--92]{Kamp1968}.
\end{proof}
\end{annexitem}
